\section{Logical meaning, verification, and reproduction}\label{sec:logic}

\subsection{The equivalence proved in ZFC}

The dependency chain is
\[
 \text{integer solution}
 \ \longleftrightarrow\ \text{finite accepted blank tableau}
 \ \longleftrightarrow\ \operatorname{Halt}(K215,\mathrm{blank})
 \ \longleftrightarrow\ \ZFCRef.
\]
The first equivalence consists of the root reconstruction, the signed-domain
arguments, the inherited Pell compiler, and the finite coefficient adapter.
Completeness runs through the chain in the opposite direction. All integer
exponents in the coefficient recipe denote fixed integers; none is a hidden
exponentiation in an unknown. Each displayed final object is consequently
an ordinary polynomial, or a finite system of ordinary polynomials, over
$\Z$.

The final endpoint is the theorem
\begin{quote}\small\ttfamily
blank\_halts\_iff\_zfc\_refutation :\\
BlankHalts machine215 $\leftrightarrow$ BProv ZFCax\_s [] Form.fBot
\end{quote}
in \src{KEndpoint215.lean}. It has no consistency hypothesis or remaining
machine-specific theorem parameter. The Rocq development gives the
corresponding equivalence with its \src{StandardZFC} theory~\cite{endpoint}.

For precision, \src{BProv} means existence of a finite list of background
axioms and a finitary first-order derivation from them. The letter B refers
to the background theory; it does not impose a numerical bound on proof
length. The namespace \src{BoundedZFCConsistency} in some imported modules
does not change this endpoint into a bounded-consistency statement.
The axiom predicate includes all instances of Separation and Replacement
and a standard selector formulation of Choice. We may define
$\Con(\mathrm{ZFC})$ using this exact effective presentation. The usual
finite proof translations between standard equivalent presentations give
the same result for another conventional coding~\cite{endpoint,t4}.

The machine calibration uses finitary syntax, first-order completeness,
constructible-universe arguments, and a well-founded quotient-model
construction. These are ordinary mathematical constructions available in
ZFC; they do not require assuming $\Con(\mathrm{ZFC})$. The new algebraic
steps use finite extension fields, automorphisms of a finite normal closure,
elementary Pell arithmetic, and finite computation. They too are
formalizable in ZFC. Composing those proofs gives Theorem~\ref{thm:main},
and classical negation gives its unsolvability version.

This formulation deliberately separates two claims. We supply an ordinary
mathematical proof of a theorem provable in ZFC. We do not supply a new
machine-checked derivation in an encoded ZFC proof calculus, or a new
end-to-end Lean/Rocq formalization of the Diophantine reductions. The
external proof-assistant endpoint by itself is not such a translation.
No stronger internalization in PA is asserted here.

If ZFC is consistent, the displayed polynomial has no solution in the
standard positive integers, or in the standard integers for its signed
variants. Under the usual hypotheses of the second incompleteness theorem,
ZFC then cannot prove this lack of solutions: the proved equivalence would
turn such a proof into a proof of its own consistency. The construction
does not determine whether a solution exists. Nor do we infer from mere
consistency that ZFC cannot prove the opposite assertion; that stronger
independence conclusion would require an additional argument or hypothesis.

\subsection{What was checked, and what was imported}

The source snapshot is pinned to the full Git commit in the abstract.
Files in \src{upstream/} preserve the inherited texts and their original
historical snapshot notices. \src{SOURCE_MANIFEST.json} records their
retrieval paths, Git blob identities, and SHA-256 hashes. The literal
430-transition machine table is included and compared with the exported
transition list before the finite compiler is regenerated.

The archive also preserves two \emph{inherited} endpoint audit logs.
The Lean log reports only \src{propext}, \src{Classical.choice}, and
\src{Quot.sound}; the Rocq log reports closure under the global context.
We inspected these records and the endpoint source. We did not execute a
fresh Lean or Rocq build in preparing this archive. Their limited size
does not make them a substitute for the complete upstream developments.

The new executable checks are standard-library Python. They verify:
\begin{enumerate}
\item Circuit register discipline, liveness, exact operation counts,
  consistency with independently evaluated norm formulas, and specialized
  polynomial degree certificates. Structural leading-form arguments in
  the proofs establish the uniform exact degrees for the authentic fixed
  coefficient tuple.
\item The signed substitutions, the negative-parameter Pell projection
  identities, the population identity used to make $s$ odd, and explicit
  small negative-scale Pell tuples with both main-root signs excluded.
\item The 26-integer norm by a separate implementation of multiplication
  and conjugation in $\Z[r,u]/(r^2-A,u^2-S)$. This checks the expanded
  circuit without copying its scalar norm formulas, including $T=0$,
  $X=0$, and negative-$s$ cases.
\item Every candidate successor for the refined admissible horizontal
  triples, the exact 5263-tile enumeration, clause/copy geometry, carry
  bounds, synchronization anchors, mask populations, and the modular unit
  needed to tune the actual rank. Eighteen finite control tableaux check
  927 seam and boundary windows.
\end{enumerate}

Finite tests support the arithmetic specification and catch implementation
errors. They do not establish the universal Pell lemmas, the unbounded
tableau theorem, or consistency. In particular, the synthetic exact-root
fixtures are tests of the algebraic norm identity; they are not claimed
solutions for the K215 coefficients. The Galois reconstruction and signed
zero-case arguments are proved in the text and independent audit notes,
not inferred from testing a sample of integers.

\subsection{Files and reproducible commands}

From the extracted archive root, run
\begin{quote}\ttfamily
python3 reproduce.py\\
latexmk -pdf -interaction=nonstopmode -halt-on-error report.tex
\end{quote}
The first command needs Python 3.10 or later and its standard library only.
It builds the new arithmetic packets and the finite compiler in a temporary
directory, compares them with the bundled artifacts, checks the literal
machine binding, runs the separate quotient-algebra audit, and writes a
verification receipt. It does not expand the enormous fixed coefficients.
The second command uses a standard TeX installation to rebuild this PDF.
The corresponding \src{Makefile} provides \src{verify}, \src{pdf}, and
\src{clean} targets.

\begin{center}\small
\begin{tabular}{>{\raggedright\arraybackslash}p{4.6cm}p{8.6cm}}
\toprule
Path & Purpose\\
\midrule
\src{report.tex}, \src{sections/} & Complete article sources.\\
\src{data/polynomial12.json} & Fully substituted 12-positive-variable product.\\
\src{data/system12.json} & Its three comparisons.\\
\src{data/integer26/} & 26-integer product, system, and sum of squares.\\
\src{data/integer27/} & Five lower-degree 27-integer alternatives.\\
\src{code/} & Deterministic circuit builders and independent norm audit.\\
\src{coefficients/} & Literal machine, 5263 tiles, coefficient recipe,
  finite verifier, and inherited endpoint evidence.\\
\src{COEFFICIENT_BINDING.json} & Exact mapping from each circuit's port names
  to the one complete coefficient tuple.\\
\src{proofs/}, \src{upstream/} & New proof/audit notes and inherited proof dependencies.\\
\src{comparison.csv} & Machine-readable comparison of the stated alternatives.\\
\src{SOURCE_MANIFEST.json} & Pinned source provenance.\\
\src{MANIFEST.sha256} & SHA-256 checksums for the distributed files.\\
\bottomrule
\end{tabular}
\end{center}

\subsection{Which form to use}

For the fewest positive unknowns in this reviewed family, use
\src{polynomial12.json}; its three-equation system is slightly shorter to
evaluate. For unrestricted integer unknowns, the new 26-integer equation
uses the fewest coordinates among the explicitly bound variants presented
here. If degree or arithmetic size matters more, the 27-integer ordinary
system has degree 46 and 107 operations, while its single ordinary product
has degree 172 and 116 operations. The unchanged 15-positive-variable
product remains the shortest known evaluator in this repository comparison,
at 67 operations. A degree-four target still favors the existing quartic
lift. The refined coefficient tuple can be used coherently with any version
whose complete blank-compiler interface is proved.

The new reductions therefore add points to a tradeoff table. They do not
make one equation uniformly simplest. Further useful work would be to
formalize the separating-map and denominator lemmas, reduce the large
weights while retaining their signed bounds, or compile the smaller
numerical proof calculus directly. Those directions are not treated as
results of this report.
